\hwhat{Does a kicked agent ignore a second kick until the task the first one launched is over? Every kick was timed at the model call that first read it (DQ1 ledger). The second kick's effect, as a function of spacing $\delta$, was compared with an isolated first kick under past-only matching. The synthetic validation (planted windows) came first. Replication covered 35 non-holdout periods, with native tests on NE43 (nudger off), G38 (pause gates) and G04 (human chat). No refractory window was found. A second kick read by a later call keeps 80--130\% of the effect at every $\delta$ from 1\,min to 4\,h: mentions $R=0.78$ [0.67, 0.89] (15 periods), human messages $R\approx1.4$, nudge re-fires 1.28. Only a kick read in the \emph{same} call adds little ($R=0.2$--0.4). The episode lock is absent (in-episode 0.67 vs 0.83). Holdout not run.}

\hmodel{Model 09 (renewal / refractory Hawkes) with model 03's refractory compartment. A kick read at time $t$ multiplies the agent's escape (or reply) hazard by a gain $g_c$ with a short kernel $\kappa$. A recovery function $r$ is measured from the previous effective kick. H43 claims $r=0$ inside the launched task episode (length $L$). The rivals are a habituation clock $r(\delta)$, $r\equiv1$ (state dependence only) and read-out batching (kicks read in one call count once). $R(\delta)=E_2/E_1$ uses pooled log hazard ratios.}

\hmath{\vspace{-5pt}\begin{gather*}
\lambda_{\rm esc}(t)=\lambda_0(a)\,[1+g_c\,\kappa(t-t_{\rm read})\,r(t-t_{\rm last})]\\
\text{H43: } r=\mathbb{1}[\,t-t_{\rm last}>L\,];\qquad \text{habituation: } r=1-a\,e^{-\delta/\tau}\\
E=\ln\frac{\sum\text{events}_K/\sum\text{at-risk}_K}{\sum\text{events}_C/\sum\text{at-risk}_C}\\
R(\delta)=E_2(\delta)/E_1,\qquad \delta_{1/2}:\ R(\delta_{1/2})=\tfrac12
\end{gather*}\vspace{-7pt}}

\hobs{../figures/summary_obs.pdf}{Second-kick effect relative to a first kick (lnHR ratio, 95\% CI), against the spacing of the two receiving calls. Diamonds: the second kick read in the \emph{same} call. Dotted: H43's prediction ($R=0$ inside a 20-min launched episode). Mentions (G51) dip only to ${\sim}0.8$ (pooled over 15 periods: 0.78); human messages (G04) and nudge re-fires (G51) stay at or above 1.}

\htelos{Useful as a negative design rule. If you run a swarm, you do not need to wait for an agent to finish its task before nudging or messaging it again. A kick read by a later call works about as well as the first. What wastes a kick is sending it before the agent's next model call, e.g.\ during its PAUSE: it lands in the same read as the earlier one and adds little. Merge such messages or deliver them later. The nudger's 15-min floor is already past any refractoriness. Its real limit is that a first nudge buys a glance (escape $\times1.75$), not sustained work. The physics label (refractory renewal) did not survive. What remains is the mundane fact that the context assembly is the unit of attention.}

\hbottom{Kicks leave no refractory window: a second kick read by a later model call works about as well as the first, and only kicks absorbed into the same read are wasted.}

\hresults{\footnotesize\begin{tabular}{@{}p{0.37\columnwidth}@{\ }p{0.5\columnwidth}@{\ }c@{}}
Prediction & Observed & \\\hline
\raggedright P1 window, $R_{\rm short}\le0.5$ & \raggedright mentions 0.57--0.79; humans, nudges $\ge1$ & $\times$\\
\raggedright P2 window $\approx$ episode & \raggedright $\delta_{1/2}=0$ vs $\tilde L$ = 5--56 min & $\times$\\
\raggedright P3 episode lock & \raggedright $R_{\rm in}$ 0.65--0.69 vs $R_{\rm act}$ 0.81--0.84 & $\times$\\
\raggedright P4 same-call kick $\le0.3$ & \raggedright 0.18 (reply), 0.31 (talk) & \checkmark\\
\raggedright P5 synthetic recovery & \raggedright nulls clean; mention windows only & $\sim$\\
\raggedright P6 NE43 nudge window & \raggedright no sustained $E_1$; re-fire $R$ = 1.28 & $\times$\\
\raggedright P7 invariance across NE43 & \raggedright mention $\Delta R=\pm0.16$ & \checkmark\\
\raggedright P8 G38 pause gates & \raggedright 8 re-kicked gates: untestable & n/a\\
\raggedright P9 G04 human trains & \raggedright $R(\delta\le2\,{\rm min})=1.46$ & $\times$\\
\raggedright P10 rule reliable & \raggedright no window in 98--100\% of draws & \checkmark\\
\end{tabular}}

\hobsb{../figures/synthetic_validation.pdf}{Synthetic validation at G51 scale. (a) Mention $R(\delta)$: the null stays near 1. A planted 30-min habituation clock is recovered ($\delta_{1/2}$ 15 vs 18\,min); an episode lock gives an intermediate curve. (b) The episode test separates the three worlds. The real data look like the null in both panels.}

\hcaveats{Only mention curves are identifiable at G51 size. Human messages are testable in G04 alone (5-min quiet rule, amended before outcomes), nudges marginally. The nudge first-kick effect exists for ``any activity'' (post hoc) but not for sustained work (pre-registered), so $R$ there rests on the post hoc outcome. Re-fired nudges outperforming first ones may reflect the nudger's unobserved text trigger. Estimated call starts for non-Gemini agents misclassify some same-call vs next-call pairs, pulling $R(\delta\le2)$ toward the batched value. Episodes come from minute-grid runs, not task content. 35 periods $\times$ 3 classes $\times$ 5 outcomes: replication verdicts are templated and not independent.}

\hnext{Price the $k$-th kick inside one read (per-provider timing). Test nudge facilitation against the trigger (H35). Run content-defined episodes and a content-drift outcome. Run \texttt{confirm.py}.}
